\section{Conclusion}
\label{sec:conclusion}

This work investigated probabilistic graph-based traffic forecasting under
distribution shift using the METR-LA dataset. A leakage-free experimental
pipeline was developed to compare a persistence baseline, a purely temporal
GRU, and a Graph-Temporal model combining diffusion-based message passing with
recurrent temporal modeling. In addition to point forecasting accuracy, the
analysis considered predictive uncertainty through quantile regression,
empirical coverage, and prediction interval width.

The experimental results show that incorporating the road-network graph
consistently improves forecasting accuracy over the Temporal model,
particularly at longer forecasting horizons. The graph ablations further
indicate that the weighted connectivity structure contains useful predictive
information, while increasing the diffusion depth provides only limited
additional benefits. These improvements, however, come with higher
computational cost and slightly reduced empirical coverage.

The high-congestion stress test reveals a more challenging limitation.
During congested periods, forecasting errors increase substantially and the
model responds by producing wider prediction intervals. Nevertheless,
coverage deteriorates at longer horizons, showing that increased uncertainty
does not necessarily translate into reliable uncertainty estimates under
distribution shift.

Post-hoc calibration through Conformalized Quantile Regression substantially
improves overall coverage and brings it close to the nominal 80\% level
without modifying the median forecasts or retraining the model. However, the
remaining undercoverage under high congestion shows that marginal calibration
alone is insufficient to fully address regime-specific reliability. A
congestion-aware training objective was therefore proposed as a possible
extension to explicitly increase the importance of these rare and difficult
traffic conditions during learning.

Overall, the results demonstrate the value of exploiting spatial information
for traffic forecasting while also highlighting the importance of evaluating
predictive uncertainty beyond average test performance. Reliable deployment
requires considering not only forecasting accuracy, but also calibration,
robustness to distribution shift, and the computational cost associated with
more expressive spatio-temporal models.